\section{Tests That Pin a Threshold to a Measurement}
\label{sec:tests-pin-thresholds}

A threshold with no test is a guess that compiles. The discipline that changes it is simple:
\textbf{every constant is justified by a number, and every number has a test on the other side
of it.}

\subsection{The Near-Miss Rule}

Any rule that suppresses or classifies something needs a case that comes \emph{close} and
must still behave correctly. For a fall detector the near misses are the interesting ones:

\begin{table}[H]
\centering
\begin{tabular}{llll}
\hline
\textbf{Boundary} & \textbf{Value} & \textbf{Must trigger} & \textbf{Must not trigger} \\
\hline
free fall & $<$ 400\,mg & real drops: 32--50\,mg & shaking: never below 630\,mg \\
free-fall length & $\geq$ 40\,ms & 42, 70, 80, 100, 138, 145, 176--252\,ms & dips of 10--25\,ms \\
landing & $\geq$ 1100\,mg & 1148--12823\,mg & caught in the hand: $\approx$ 990\,mg \\
stillness & $<$ 1500\,mg for 600\,ms & 1.04--1.17\,s observed & bouncing: $>$1500\,mg \\
\hline
\end{tabular}
\caption{Every constant, the measurements on both sides of it, and the control that must not fire.}
\label{tab:thresholds-near-miss}
\end{table}

Two of these margins are uncomfortably thin, and saying so is part of the test:

\begin{warningbox}
    The landing threshold is 1100\,mg with the softest real landing at 1148\,mg — an 8\,\%
    margin, on nine drops, performed by one person on one desk. It is enough to ship a rule
    and not enough to claim the rule is right. The test encodes the measurements, not a
    belief about all landings everywhere.
\end{warningbox}

\subsection{Writing the Test From the Recording, Not From the Idea}

The failing pattern to avoid is a test that restates the implementation:

\begin{minted}{rust}
// Useless: this asserts that the code does what the code does.
#[test]
fn detects_a_fall() {
    let mut det = FallDetector::new();
    assert!(det.update([0, 0, 0], 5).is_some());   // whatever the code says
}
\end{minted}

The useful version uses a sequence of samples taken from a recording, and asserts on the
number a person measured:

\begin{minted}{rust}
#[test]
fn a_soft_landing_after_a_short_free_fall_is_detected() {
    // Drop 3 of the 400 Hz round: 42 ms of free fall, landing of only 1493 mg,
    // then quiet. This is the case an 800 ms settle rule lost.
    let mut det = FallDetector::new();
    feed(&mut det, [990, 40, 60], 300);      // resting
    feed(&mut det, [30, 25, 20], 42);        // in the air
    feed(&mut det, [1200, 700, 400], 20);    // 1493 mg landing
    match feed(&mut det, [980, 40, 60], 700) {
        Some(Event::Fall { peak_mg, .. }) => {
            assert!((1200..1700).contains(&peak_mg), "peak {peak_mg}");
        }
        other => panic!("expected a fall, got {other:?}"),
    }
}
\end{minted}

Three properties make this test worth its lines: the numbers come from hardware, the
assertion range is wider than the measured value (sensor noise is real), and it fails loudly
with the actual event when the behaviour changes.

\subsection{Embed the Recording in the Image}

Host tests prove the logic. They do not prove that the code on the device behaves the same,
because the device may run a different build, a different clock, a different optimisation
level. The cheap bridge between the two is to embed the recordings in the firmware and replay
them at boot:

\begin{minted}{text}
=== self test at boot, no host attached ===
real=clean  synth=ok  200Hz=2/2  400Hz=4/4
\end{minted}

That line is the whole argument: the handling recording produces nothing, the synthetic fall
is detected, and six of the real drops — replayed through the same binary that runs live — are
all found. A regression in the shipped image shows up as a different line, on a device that
needs no laptop.

\begin{tipbox}
    Embed the recordings as \codeinline{const} arrays (a few kilobytes) and replay them
    \emph{before} the main loop. It costs one flash write and turns every power-up into a
    regression test.
\end{tipbox}
